\section{Computation normalization and the bridge to the existing core}
\label{sec:normal-bridge}

The proof above is an exact comparison on canonical sequents. We now make
explicit why that comparison retains full iPRA, including the two places
where a naive normalization argument would be incomplete.

\subsection{What the normalizer does}

Raw scalar terms are finite expressions $x$, $0$, $S(t)$ and
$I(n,a;u.h)$. The binder $u$ binds only in $h$, not in $n$ or $a$.
Reduce both computation rules in every numerical subterm:
\[
 I(0,a;u.h)\longrightarrow a,\qquad
 I(Sn,a;u.h)\longrightarrow h[I(n,a;u.h)/u].
\]
Reduction is allowed inside an iterator body. A canonical term has no
iterator whose counter is headed by $0$ or $S$. Formula normalization
changes terms and preserves equality, $Q$ and $D$ roots.
Normalizing a clause treats its consequent under its own bound variable.

The full-normalization package supplies the detailed termination and
confluence argument \cite{normal}; the following proof records its
ingredients used here.

\begin{lemma}[Termination, confluence, and substitution coherence]
Every well-scoped scalar term has a unique computation normal form.
For capture-free substitution $\sigma$,
\[
 \nf(t\sigma)=\nf\bigl(\nf(t)\,\nf(\sigma)\bigr).
\]
The same statement holds for formulas and clauses with their binders respected.
\end{lemma}
\begin{proof}
Use a computability argument for ground-type iteration.
A body $h(u)$ is computable when substituting any computable numerical
term for $u$ produces a strongly normalizing term; passive parameters
are handled by a computable valuation. Variables, zero and successor
preserve computability.

For iterator closure, suppose $n,a$ are strongly normalizing and
$h$ has that substitution property. Let $\nu(t)$ be the maximum
reduction length of a strongly normalizing term. Induct lexicographically on
\[
 \bigl(\nu(n),\,|n|,\,\nu(a)+\nu(h)\bigr).
\]
A reduction inside the counter decreases the first coordinate.
A reduction inside the seed or body decreases the last; the computability
property of a reduced body follows by substitution compatibility of reduction.
The zero head reduct is $a$. For a successor counter $Sm$, the smaller
counter $m$ has the same reduction height and smaller syntax size.
The induction hypothesis makes $I(m,a;u.h)$ computable, and the body's
substitution property makes the head reduct
$h[I(m,a;u.h)/u]$ strongly normalizing.
Structural induction on terms under computable valuations completes the
argument. Free variables themselves are computable, so open terms are covered.

The two root patterns are disjoint and left linear. Reductions in disjoint
subterms commute. A reduction inside a substituted body or argument joins
the corresponding duplicated or erased reductions after a head step.
This gives local confluence; termination gives confluence.
Finally substitution takes a computation reduction to a finite reduction
sequence, possibly empty when an erased part is affected. Both sides of the
displayed identity reduce from substitution instances of the same raw term,
so their unique normal forms agree.
\end{proof}

No theorem that the \emph{uniform} normalizer is primitive recursive is
assumed. Every fixed iterator description represents a primitive recursive
operation, whereas normalization takes arbitrary descriptions as input.
The total normalizer is part of the trusted algorithm and is counted in the
source measurements. A practical run may exhaust time or memory on a large
certificate. The mathematical calculus permits every finite derivation;
the Python programs implement it subject to ordinary resource limits.

\subsection{Restoring every case introduction}

The single introduction I selects progress for an equation at the literal
zero witness. A raw witness can become zero under normalization.
It is therefore incorrect simply to normalize a raw I node and assume that
its mode stayed unchanged.

First work in $\CT$, where crossed transport is available. For an equation
$E=(s=t)$ put
\[
 G(v,E):=(C(v,0,s)=C(v,0,t)),\qquad
 L(E):=q_vG(v,E)=D_{a,v}G(v,E),
\]
with unused $a$ and fresh binders. The zero instance of $G$ is literally
$0=0$; its successor instance is literally $E$. Hence
\[
 E\vdash L(E),\qquad L(E)\vdash E.
\]
The first direction uses the equality-free case-closure interface at dummy
witness $1$; the second eliminates that closure at $1$ using reflexivity.
These conversions use neither equation induction nor equality transport.

Suppose the desired missing case has premise
$\Gamma,E(0,0)\vdash E(0,Sz)$. Wrap its base and conclusion in $L$.
The matrix $L(E(x,y))$ has a compound $D$ root, so I at witness zero uses
cases and gives $D_{x,y}L(E(x,y))$. Eliminate this at a fresh variable $b$.
To introduce $D_{x,y}E(x,y)$ at witness $b$, turn its assumed base
$E(b,0)$ into $L(E(b,0))$, use the retained clause at $Sz$, and unwrap
the result. Witness $b$ is not literally zero, so the final I is again
in case mode. Close the elimination.

Thus every case introduction is derivable. Progress for an arbitrary
equation $E(y)$ uses I with an unused witness coordinate and witness zero,
then the $q$-elimination interface. This supplies contextual equation
induction once a base is given.

\subsection{An auxiliary calculus removes mode ambiguity}

Use an auxiliary canonical calculus with separate unrestricted case
introduction and equation-progress introduction. Its terms and other
rules are unchanged. The preceding constructions translate it into
$\CT$. Conversely a case-mode I is unrestricted case introduction, and
a progress-mode I is equation progress followed by case introduction.
Thus the two canonical calculi have the same consequences.

Translate the raw native core into the auxiliary canonical calculus:
normalize numerical terms, replace $Q_yB$ by $q_yB$ in the clausal
presentation, and translate ordinary assumptions to constant clauses.
Computation axioms become reflexivities. A raw transport instance
becomes a canonical one by normalized-substitution coherence.
Closure elimination becomes $q$-elimination.
An eliminated $D$ supplies its retained clause, from which the corresponding
$q$ assumption can be reintroduced. The case and progress rules are handled
separately, so no witness-mode invariance is presumed.

For the reverse comparison, every reduction represents a provable
equation in the raw core. Reducing inside a seed or counter uses ordinary
congruence. Reducing inside a \emph{bound update body} requires the following
derived fact, rather than an uncharged function-extensionality axiom.

\begin{lemma}[Iterator-body extensionality in the raw theory]
If the raw theory proves $h(u)=k(u)$ with $u$ fresh for the retained
context and passive parameters, it proves
\[
 I(n,a;u.h)=I(n,a;u.k).
\]
\end{lemma}
\begin{proof}
Apply literal-equation induction to equality of the two trajectories.
Their initial values are both $a$. At the successor step use their
computation equations, the induction hypothesis, and the proved pointwise
body equality. All passive parameters stay fixed.
\end{proof}

This proves $t=\nf(t)$ for each raw term, including reductions inside
update bodies, and lifts to formula and clause equivalence. It completes the
raw/canonical comparison. The inherited written source arguments are
included under \code{reference/}.

\begin{theorem}[Faithful comparison with the native core]
The native scalar core and $\CT$ translate proofs in both directions.
Both formula composites are internally interderivable. The numerical
maps are normalization and raw inclusion, whose composites are
internally equal to the identity. The same holds for
$\CV$ and $\CZ$ by their exact equality-rule comparisons.
\end{theorem}

\section{The two interpretations with full first-order iPRA}
\label{sec:interpretations}

\subsection{Ordinary realization}

Let $T_I$ be intuitionistic first-order logic with actual equality over
the scalar iterator language, both computation equations, contextual
induction for literal equations, coverage and separation:
\[
 \forall x\,(x=0\lor\exists y\,x=Sy),\qquad
 (0=S0)\to\bot.
\]
The separation sentence is explicit. Positive equations alone admit a
singleton model, so the ordinary falsity bridge cannot be silently omitted.

The formula and clause readings are
\[
 (s=t)^\circ=(s=t),\quad
 (D_{x,y}B)^\circ=\exists x\bigl(B^\circ(x,0)\to\forall y\,B^\circ(x,y)\bigr),
\]
\[
 [A;y.B]^\circ=A^\circ\to\forall y\,B^\circ(y),\qquad
 (Q_yB)^\circ=B^\circ(0)\to\forall y\,B^\circ(y).
\]
The interpretation domain is $x=x$; numerical equality is actual identity.

Every rule is sound. Refl is identity reflexivity. VSub follows from
equality replacement and symmetry. Use and E are implication, universal,
and existential elimination. Case-mode I uses coverage. Progress-mode I
uses induction for its literal equation under a temporary zero-instance
assumption. The eigenvariable conditions keep the retained context and
witness fixed. Every normalization step is already an equation of $T_I$.

\subsection{The inherited arithmetic compiler}

The pinned native development proves $T_I$ bi-interpretable with iPRA
\cite{repo,bridge}. Its arithmetic dependency order is significant:
first a scalar predecessor, then reflection of actual equality into a
numerical zero test, then quantifier-free induction, then pairing and
the primitive-recursive compiler. None is assumed to prove the new
variable-source reduction.

In iPRA interpret each fixed iterator by ordinary primitive recursion
ignoring the counter in its update:
\[
 i_h(0,a,\vec p)=a,\qquad
 i_h(Sn,a,\vec p)=h(i_h(n,a,\vec p),\vec p).
\]
Each fixed finite body has its own primitive-recursive description.
There is no universal total evaluator on arbitrary description codes.

In the reverse direction the inherited scalar construction supplies
pairing and projection terms $\pi,L,H$, with
$L(\pi(a,b))=a$ and $H(\pi(a,b))=b$ proved internally.
For a fixed primitive-recursive definition
\[
 f(0,\vec p)=g(\vec p),\qquad
 f(Sn,\vec p)=h(n,f(n,\vec p),\vec p),
\]
define
\[
\begin{split}
 V(n,\vec p)&=I\bigl(n,\pi(0,g^*(\vec p));
       u.\pi(SL(u),h^*(L(u),H(u),\vec p))\bigr),\\
 f^*(n,\vec p)&=H(V(n,\vec p)).
\end{split}
\]
Equation induction proves $L(V(n,\vec p))=n$; the projection identities
give the recursion equations. Zero, successor, projections and composition
translate directly. The construction recurses externally on a finite
definition. It does not replace the numerical domain by a coding range.

\subsection{Full intuitionistic logic and the composites}

The native development defines formula selectors $\operatorname{sel}(v,A,B)$
and gates $\operatorname{gate}(v,A)$ by structural recursion on formulas.
At zero the selector is $A$, and at a successor it is $B$; the gate is
true at zero and equivalent to $A$ at a successor. For equations the
definitions apply $C$ to the two endpoints. Compound formulas are padded
and their binders aligned before recursing. This is a finite syntactic
construction, not quantification over predicates.

With $q_yB=D_{w,y}B$ for an unused $w$, the logical macros are
\[
\begin{aligned}
 \bot_*&=(0=S0),\\
 A\Rightarrow_*B&=q_y\operatorname{sel}(y,A,B),\\
 A\wedge_*B&=q_y\operatorname{gate}
          \bigl(y,\operatorname{sel}(P(y),A,B)\bigr),\\
 \exists_*x\,A(x)&=D_{x,y}\operatorname{gate}(y,A(x)),\\
 A\vee_*B&=\exists_*v\,\operatorname{sel}(v,A,B),\\
 \forall_*x\,A(x)&=A(0)\wedge_*q_xA(x).
\end{aligned}
\]
Dummy coordinates are fresh, and substitutions avoid capture.
The inherited derivations in \code{reference/inherited_choice_basis.tex}
provide their intuitionistic natural-deduction rules. Our exact sequent
comparisons preserve those derivations.

For the numerical round trips, two translations of a fixed operation
satisfy the same initial and recursion equations. Equation induction
identifies their trajectories. External induction on a finite operation
description proves that its double translation equals the original
operation on \emph{every internal input}. The argument also accounts for
changes of staged body descriptions after substitution.
Formula structural induction then proves both logical round trips.

\begin{theorem}[Bi-interpretability and equiconsistency]
The ordinary realizations of $\CV$, $\CZ$, and their native six-rule
companions are bi-interpretable with full first-order iPRA.
Both numerical interpretation domains are whole, equality is actual
equality, and in both theories $x=y$ defines the numerical composite
isomorphism. Contradiction translates to $0=S0$, so the presentations
are equiconsistent.
\end{theorem}
\begin{proof}
Compose the exact comparisons with the explicit native arithmetic and
logical interpretations. Normalization is internally equal to raw
inclusion in the reverse direction, and the equality-rule changes are
the identity on end sequents. The inherited numerical and formula
round trips remain valid. The identity numerical relation is internally
total, bijective, and preserves every interpreted operation.
Falsity is preserved in both directions.
\end{proof}

This is a written mathematical proof extending the pinned formalized native
bridge. The new restricted kernels and proof translations have not been
checked by Lean or Rocq.

